\section{Method and Evaluation}\label{sec:attack:method}

The attack trains one affine map; two later comparisons train a two-layer map (Section~\ref{sec:attack:negative}) and a projector for a further decoder (Section~\ref{sec:attack:leak}).
The map sends each hidden state of the target model into the hub; the frozen projector of a decoder maps the result into that decoder's input space, and the frozen decoder reconstructs the prompt: Pythia block \vf{12} $\to$ affine map $\to$ hub $\to$ projector $\to$ decoder.

\paragraph{Map.} The map sends a hidden state $s$ of the target model to a point $\hat m$ of the hub:
\begin{equation}
\hat m = W(s - \mu) + b ,
\label{eq:attack:map}
\end{equation}
where $\mu$ is the mean hidden state of the unpaired pool, $W$ is a matrix from \vf{2048} to \vf{512} dimensions, and $b$ is a bias. % src: src/v7/adapter.py; src/v7_2/preprocess.py (centering on pool U); configs/v7.yaml (width 2048, hub_dim 512)
Only $W$ and $b$ are trained.
The map acts on one hidden state at a time, as the hub does (Section~\ref{sec:hub:design}).

\paragraph{Loss.} The loss combines a supervised term on the rows of the leak with a matching term on the unpaired pools:
\begin{equation}
\mathcal{L}(\alpha) = \alpha\,\big(\mathcal{L}_{\text{sup}} + \lambda\,\mathcal{R}\big) + (1-\alpha)\,\mathcal{L}_{\text{match}} .
\label{eq:attack:loss}
\end{equation}
The supervision fraction $\alpha$ weights the two terms. % src: results/v7_2/main/mixed-v7_9-s5.json (definition)
At $\alpha = 1$ the loss reads the leak only, and the result is called the supervised fit.
At $0 < \alpha < 1$ the loss reads the leak and the unpaired pools, and the result is called the hybrid fit.
At $\alpha = 0$ the loss reads no row, and the result is called a blind map.

\paragraph{Supervised term.} $\mathcal{L}_{\text{sup}}$ is the mean squared distance between the map's output and the hub output of the same row, divided by $\rho_H^2$.
Rows are sampled by prompt first and then within the prompt, so each prompt of the leak has equal weight whatever its number of rows. % src: src/v7_2/mixed.py (SEED_RULE: prompt uniform, then row uniform)
The constant $\rho_H$ is the root-mean-square distance of the pooled hub outputs from their mean, so the term does not depend on the scale of the hub. % src: src/v7_2/context.py (rho_H); src/v7/loss.py
The ridge penalty $\mathcal{R}$ is the squared Frobenius norm of $W$ divided by $\rho_H^2$, with $\lambda = 1$. % src: configs/v7_9_s4.yaml (mixed.ridge_lambdas); src/v7_2/mixed.py
Unless a table states otherwise, the supervised fit includes the ridge penalty.

\paragraph{Matching term.} $\mathcal{L}_{\text{match}}$ reads a batch of $B$ mapped hidden states $\hat m_1 \ldots \hat m_B$ from the unpaired pool of the target model and a batch of $B$ hub outputs $h_1 \ldots h_B$ from the unpaired pool of the hub.
No row of the first batch is matched to a row of the second.
The score of a pair is its negative squared distance, $\sigma(u, v) = -\lVert u - v\rVert^2 / \rho_H^2$.
Each mapped state is first replaced by a weighted average of the hub batch, with weights that favor the nearest hub outputs:
\begin{equation}
\bar h_i = \sum_{j} p_{ij}\, h_j ,
\qquad
p_{ij} = \frac{\exp(\sigma(\hat m_i, h_j)/\tau)}{\sum_{k}\exp(\sigma(\hat m_i, h_k)/\tau)} .
\label{eq:attack:match}
\end{equation}
The loss then asks that the average $\bar h_i$ identify the mapped state it came from among the $B$ mapped states of the batch:
\begin{equation}
\ell_i = -\log\frac{\exp(\sigma(\bar h_i, \hat m_i)/\tau)}{\sum_{k}\exp(\sigma(\bar h_i, \hat m_k)/\tau)} .
\label{eq:attack:return}
\end{equation}
$\mathcal{L}_{\text{match}}$ averages $\ell_i$ over the batch and adds the same loss with the two batches exchanged, each with weight one half. % src: src/v7_2/abc_sqdist.py (abc_sqdist_loss)
The matching term is low when mapped states and hub outputs are close to one another in a consistent way; it reads no text and no row.

\paragraph{Temperature.} The scale of the temperature $\tau$ comes from the hub alone, and its multiplier is selected by word-F1 (Section~\ref{sec:attack:hybrid}).
Let $d$ be the median gap between the highest and the second-highest score of a hub output against the other hub outputs of its batch, over \vf{100} batches of \vf{512} hub outputs. % src: src/v7_2/gap.py; results/*/v7_2/main/mixed-v7_9-s0.json gap
The gap is \vf{0.0576} on Corpus A and \vf{0.0648} on legal128w, and $\tau$ is a multiple of $d$. % src: results/v7_2/main/mixed-v7_9-s0.json gap.sqdist.median (0.05761); results/legal128w/v7_2/main/mixed-v7_9-s0.json (0.06485)

\paragraph{Training.} Every fit runs \vf{3{,}000} AdamW updates with a learning rate that rises to $3\times10^{-4}$ over \vf{300} warmup updates and follows a cosine decay to zero, no weight decay, and gradients clipped at norm \vf{1.0}. % src: configs/v7_9_s4.yaml (training); src/v7/training.py (lr_multiplier)
Each update samples \vf{512} rows of the leak, with replacement, and, for the matching term, $B$ rows from each unpaired pool. % src: results/v7_2/main/mixed-v7_9-s5.json setting (paired_batch_rows 512)
The map after the last update is evaluated; no checkpoint is selected.
A seed fixes the initial map, the draw of the leak, and the order of the batches.

\paragraph{Settings.} Two settings of the hybrid fit appear in the chapter.
On Corpus A, Sections~\ref{sec:attack:budget} to~\ref{sec:attack:hybrid} report $\alpha = 0.95$, $\tau = 2d$, and $B = 2{,}048$, selected step by step in Section~\ref{sec:attack:hybrid}. % src: results/v7_2/main/mixed-v7_9-s5.json setting, selection
Section~\ref{sec:attack:leak} reports $\alpha = 0.99$, $\tau = d$, and $B = 8{,}192$ on both corpora. % src: results/legal128w/v7_2/main/mixed-oct_w5-s1-L90.json setting; results/v7_2/main/mixed-oct_w5b-A-prefix.json setting
The result files do not record how the second setting was chosen. \todo{Author: state the selection of the second setting, or keep the sentence.}

\paragraph{Data.} The corpora are Corpus A and legal128w of Chapter~\ref{ch:hub}, each with its own frozen hub.
All training data come from the \vf{40{,}000} training prompts of a corpus. % src: results/v7_2/main/manifest.json (training_pools); results/v7_2/main/anchors-linear.json (pool_size 40000)
The training prompts are divided once, at random, into two halves of \vf{20{,}000}: one half supplies the unpaired pool of the target model, as hidden states, and the other half supplies the unpaired pool of the hub, as hub outputs of OLMo.
A leak draws its paired sequences at random from all \vf{40{,}000} training prompts, and the leaks of one seed are nested: a larger budget contains the smaller ones. % src: src/v7_2/anchors_linear.py (nested permutation); mixed-v7_9-s5.json anchor_counts (prompts_in_U)
On Corpus A, ten paired sequences hold \vf{92} rows. % src: results/v7_2/main/mixed-oct_w5b-A.json anchor_counts (seed 1)

\paragraph{Decoding and score.} Evaluation follows Section~\ref{sec:hub:eval}: each of the three decoders generates greedily from the projected output of the map, without the ground-truth prefix, up to the memory length plus \vf{8} tokens on Corpus A and plus \vf{32} tokens on legal128w. % src: configs/v7.yaml (cap_margin 8); configs/v7_legal128w.yaml (cap_margin 32)
The score is word-F1, the word overlap between the reconstruction and the true prompt defined in Section~\ref{sec:hub:eval}, averaged over prompts and over the three decoders. % src: src/v7_2/f1.py; src/v7_2/factorial.py
In this chapter, word-F1 without a decoder's name denotes that three-decoder mean.

\paragraph{Control.} The mismatched-prompt control gives each decoder the mapped states of another prompt and scores the output against the original prompt.
The other prompt is the evaluation prompt whose number of rows is nearest; Chapter~\ref{ch:hub} takes the next prompt in order instead. % src: src/v6_4/controls.py (length_matched: nearest length, lower index); mixed-*.json cells.*.wrong_memory.rule
A map whose word-F1 equals its control gains nothing, in mean word-F1, from the prompt's own mapped states.

\paragraph{Hub reference.} The hub reference is the word-F1 that the three decoders reach from the hub outputs of OLMo for the evaluation prompts.
It measures the frozen hub without any map and averages the three directions from OLMo, including OLMo to itself, where Chapter~\ref{ch:hub} averages the six cross-model directions.
It is a reference, not an upper bound: the target model and OLMo tokenize a prompt differently, so even a map that fitted every row exactly would give the decoders a different sequence.
The hub reference is \vf{0.946} on Corpus A and \vf{0.982} on legal128w. % src: results/v7_2/main/factorial.json controls.true_olmo_code.macro_f1 (0.9459); results/legal128w/v7_2/main/mixed-v7_9-s0.json setting.corpus.g2.macro_f1 (0.9824)

\paragraph{Prompt sets.} Two disjoint sets of validation prompts are read, and neither overlaps the training prompts.
The selection prompts are \vf{500} prompts drawn at random from the \vf{5{,}000} validation prompts of a corpus. % src: v7_run.py (subset_seed 20260903, dev_prompts 500); configs/v7.yaml; configs/v7_legal128w.yaml
They selected the settings of Section~\ref{sec:attack:hybrid}, and they also report every score in this chapter unless the text names the held-out prompts.
The held-out prompts are the other \vf{4{,}500} validation prompts. % src: src/v7_2/context.py; results/legal128w/v7_2/main/eval/confirmation/prompt_ids.json
They selected no setting of the attack, although all \vf{5{,}000} validation prompts selected the hub checkpoint in Chapter~\ref{ch:hub}; they report the legal128w leaks of Section~\ref{sec:attack:leak} for which the text gives a held-out score.
A score on the selection prompts is therefore not held out from selection.
The test split of each corpus has not been read.

\paragraph{Intervals and seeds.} A difference between two fits is reported with a \vf{95} percent percentile interval from a paired bootstrap over prompts, with \vf{1{,}000} resamples on the selection prompts and \vf{2{,}000} on the held-out prompts. % src: mixed-*.json bootstrap (1000); results/legal128w/v7_2/main/oct_w5-summary-confirmation-A.json bootstrap (2000)
The intervals describe variation over prompts, not over seeds.
Sections~\ref{sec:attack:budget} and~\ref{sec:attack:hybrid} report three seeds; Section~\ref{sec:attack:leak} reports one seed, and each table states its seeds.
